\documentclass[12pt,a4paper]{ctexart}

\usepackage[margin=1in]{geometry}
\usepackage{setspace}
\usepackage{hyperref}
\usepackage{enumitem}
\usepackage{xcolor}

\hypersetup{
  colorlinks=true,
  linkcolor=blue,
  urlcolor=blue
}

\setstretch{1.25}

\title{项目中期报告\\基于大语言模型的智能游戏代练服务平台}
\author{徐翱 \quad 易熙铭}
\date{}

\begin{document}
\maketitle
\tableofcontents
\newpage

\section{项目目标与技术选型}

\subsection{项目目标}
本项目延续初期报告提出的总体目标，围绕“需求理解、订单流转、角色审核、后台治理”四条主线推进实现。中期阶段的重点不是继续扩张功能边界，而是先完成一个可演示、可部署、可继续迭代的核心业务闭环。结合当前代码仓库，本阶段主要聚焦以下内容：

\begin{itemize}[leftmargin=2em]
  \item 完成用户注册、登录、个人资料维护与密码修改等基础账户能力；
  \item 接入大语言模型，对用户自然语言描述的代练需求进行结构化解析；
  \item 实现订单创建、查询、接单、完单、取消与争议处理等核心业务流程；
  \item 补齐普通用户申请成为代练、管理员审核、代练名额分配等角色流转机制；
  \item 提供基于 Docker Compose 的开发与部署方式，保证项目具备较好的工程可复现性。
\end{itemize}

\subsection{技术选型}
结合课程项目的时间约束和实现成本，当前版本采用了较为稳妥的前后端分离方案：

\begin{itemize}[leftmargin=2em]
  \item 后端：FastAPI + SQLAlchemy + Pydantic + Alembic + MySQL。该组合具备接口定义清晰、类型约束明确、迁移机制完善等优点，适合快速搭建业务型系统；
  \item 前端：Vue 3 + Vite + Pinia + Vue Router + Axios + TailwindCSS，用于构建轻量、可维护的单页应用；
  \item AI 能力：通过 OpenAI 兼容 SDK 接入 DeepSeek 接口，用于提取游戏名称、当前段位、目标段位、预算等结构化信息；
  \item 部署方式：Docker + Docker Compose，同时保留生产模式与开发模式两套编排配置。
\end{itemize}

\section{后端代码结构与实现}

\subsection{代码分层}
当前仓库中的后端采用较为清晰的分层结构，各层职责边界明确，便于后续维护与扩展：

\begin{itemize}[leftmargin=2em]
  \item \texttt{app/api}：负责路由定义、依赖注入与鉴权控制，对外暴露 REST API；
  \item \texttt{app/services}：封装业务逻辑，包括用户服务、订单服务与 AI 解析服务；
  \item \texttt{app/models}：定义数据库 ORM 模型，当前核心实体为 \texttt{User} 与 \texttt{Order}；
  \item \texttt{app/schemas}：定义请求与响应的数据模型，统一输入输出约束；
  \item \texttt{app/core} 与 \texttt{app/db}：负责配置读取、安全能力、数据库会话与应用启动配置。
\end{itemize}

\subsection{核心模块}
从当前实现看，后端已经形成了较完整的业务骨架：

\begin{itemize}[leftmargin=2em]
  \item 认证模块：支持注册、登录、刷新令牌、获取当前用户、更新个人资料与修改密码，并在应用启动时补齐默认管理员账户；
  \item 订单模块：支持需求分析、创建订单、分页查询、查看详情、修改订单、接单、完单、取消和争议处理；
  \item 用户模块：支持普通用户提交代练申请，上传段位截图，并查询自身申请状态；
  \item 管理模块：支持管理员审核代练申请、分配代练名额，以及对异常订单执行取消或争议干预；
  \item AI 模块：调用 DeepSeek 对中文需求文本进行结构化提取，同时对违规内容进行初步拦截。
\end{itemize}

\subsection{迁移与工程化支撑}
数据库结构通过 Alembic 管理，目前至少包含以下两个关键版本：

\begin{itemize}[leftmargin=2em]
  \item \texttt{001\_initial}：初始化用户表与订单表；
  \item \texttt{002\_booster\_application}：补充代练申请、审核记录、代练名额等字段与关联关系。
\end{itemize}

除业务接口外，后端还补充了若干工程化能力，例如：应用启动时自动创建上传目录、挂载 \texttt{/uploads} 静态资源路径、提供健康检查接口，以及将常见参数校验错误转换为中文提示。这些工作虽然不直接体现为业务功能，但对部署稳定性和演示体验有明显帮助。

\section{前端代码结构与实现}

\subsection{页面与路由}
前端目前已经围绕核心业务闭环搭建了对应页面，主要包括：

\begin{itemize}[leftmargin=2em]
  \item 首页、登录页、注册页；
  \item 订单列表页、订单创建页、订单详情页；
  \item 个人中心页面，用于资料维护与代练申请提交；
  \item 管理员页面，用于审核代练申请与处理订单；
  \item 兜底的 \texttt{NotFound} 页面，用于处理非法路由访问。
\end{itemize}

\subsection{状态管理与网络层}
当前前端主要使用 Pinia 管理状态，并通过 Axios 统一处理接口请求：

\begin{itemize}[leftmargin=2em]
  \item \texttt{auth} Store 负责维护用户信息、访问令牌与刷新令牌，并对登录态进行初始化；
  \item \texttt{orders} Store 负责订单列表、订单详情、分页信息与 AI 分析结果的状态管理；
  \item Vue Router 在路由守卫中统一处理登录校验、管理员权限校验与页面标题切换；
  \item Axios 拦截器会自动注入 JWT，并在遇到 \texttt{401} 时尝试刷新令牌，失败后再回退至登录页。
\end{itemize}

\subsection{前后端联动流程}
结合当前页面与接口实现，可以归纳出以下典型业务链路：

\begin{enumerate}[leftmargin=2em]
  \item 普通用户注册并登录后，可在“发布订单”页面先输入自然语言需求，再由 AI 生成可编辑的结构化表单；
  \item 用户确认表单后创建订单，并在订单列表与订单详情页面查看状态变化；
  \item 用户可在个人中心提交代练申请，上传段位截图，等待管理员审核；
  \item 管理员审核通过后，用户角色变更为代练，并获得可接单名额；
  \item 代练可以浏览待接订单、执行接单与完单操作；当订单出现异常时，用户或管理员可继续触发取消或争议流程。
\end{enumerate}

\section{数据库设计（核心）}

\subsection{\texttt{users} 表}
\texttt{users} 表承担了账户信息与角色流转两类职责。除邮箱、用户名、密码摘要、角色、启用状态等基础字段外，本阶段新增了一组面向代练审核流程的字段，包括：

\begin{itemize}[leftmargin=2em]
  \item 申请状态：\texttt{booster\_application\_status}；
  \item 申请材料：\texttt{booster\_application\_game}、\texttt{booster\_application\_current\_rank}、\texttt{booster\_application\_target\_rank}、\texttt{booster\_application\_proof\_url}、\texttt{booster\_application\_note}；
  \item 审核记录：\texttt{reviewed\_by\_admin\_id}、\texttt{reviewed\_at}、\texttt{review\_note}；
  \item 名额控制：\texttt{booster\_quota}。
\end{itemize}

这样的设计将“用户身份”和“身份演进过程”统一收敛在用户实体中，便于管理员审核、角色切换与权限控制。

\subsection{\texttt{orders} 表}
\texttt{orders} 表负责描述代练订单的完整生命周期。当前实现中的关键字段主要包括：

\begin{itemize}[leftmargin=2em]
  \item 关联字段：\texttt{user\_id}、\texttt{booster\_id}；
  \item 业务字段：\texttt{game\_name}、\texttt{current\_rank}、\texttt{target\_rank}、\texttt{price}、\texttt{status}；
  \item 描述字段：\texttt{description\_raw}、\texttt{description\_ai}、\texttt{notes}；
  \item 时间字段：\texttt{created\_at}、\texttt{updated\_at}、\texttt{locked\_at}、\texttt{completed\_at}；
  \item 敏感字段：\texttt{game\_account}、\texttt{game\_password}，其中游戏密码在服务层写入前会进行加密处理。
\end{itemize}

\subsection{枚举设计}
为了降低魔法字符串带来的维护成本，当前版本将多个核心状态抽象为枚举类型：

\begin{itemize}[leftmargin=2em]
  \item 用户角色：\texttt{USER / BOOSTER / ADMIN}；
  \item 订单状态：\texttt{PENDING / LOCKED / COMPLETED / DISPUTED / CANCELLED}；
  \item 代练申请状态：\texttt{NONE / PENDING / APPROVED / REJECTED}。
\end{itemize}

这种方式不仅提升了代码可读性，也为后续的权限判断、状态机约束和接口文档生成提供了统一基础。

\section{本阶段已完成的优化}

\begin{itemize}[leftmargin=2em]
  \item 权限收敛：公开注册接口即使接收到高权限角色参数，也会统一降级为 \texttt{USER}，避免通过前端绕过角色创建规则；
  \item 敏感信息处理：账户密码采用哈希存储，游戏密码在订单写入与更新时进行加密，降低明文泄露风险；
  \item 并发一致性：代练接单逻辑使用数据库行级锁，避免同一订单被重复接取；
  \item 业务约束补强：接单时同时校验角色、订单状态、代练名额与“不能接自己的订单”等边界条件；
  \item 审核闭环完善：新增代练申请、管理员审核与名额分配流程，替代早期依赖手工修改角色的做法；
  \item 配置外置化：后端 CORS 白名单、前端 API 基址、默认管理员信息等均通过环境变量管理；
  \item 容器兼容性优化：前端镜像构建兼容存在或不存在 \texttt{package-lock.json} 的两种场景，后端启动时主动创建上传目录，降低容器冷启动失败概率；
  \item 交互体验优化：补充中文化校验提示、管理员面板与个人中心申请入口，使演示路径更加完整。
\end{itemize}

\section{遇到的问题与排障过程}

\subsection{构建与配置问题}
\begin{itemize}[leftmargin=2em]
  \item 前端镜像构建阶段曾遇到 \texttt{npm ci} 与无锁文件仓库不兼容的问题，后续在 Dockerfile 中改为“有锁文件用 \texttt{npm ci}，无锁文件回退 \texttt{npm install}”；
  \item 后端调试参数若直接复用通用环境变量 \texttt{DEBUG}，容易受到宿主机环境影响，因此部署层改为通过 \texttt{BACKEND\_DEBUG} 间接注入，减少误配置风险。
\end{itemize}

\subsection{运行时问题}
\begin{itemize}[leftmargin=2em]
  \item 后端需要挂载上传目录为静态资源路径，若目录在容器启动时不存在，会影响应用初始化；当前版本已在启动阶段主动创建 \texttt{uploads} 目录；
  \item 数据库迁移在 MySQL 环境下需要额外注意 DDL 的非事务特性，若迁移中断，可能出现“结构已部分执行但版本号未同步”的状态，因此需要在排障时同时核对表结构与 Alembic 版本信息。
\end{itemize}

\subsection{业务联调问题}
\begin{itemize}[leftmargin=2em]
  \item 代练无法接单的问题并非单点故障，而是角色、状态、名额和自接单限制等多重条件共同作用的结果；本阶段已将这些约束统一收敛到服务层处理；
  \item 用户、代练与管理员三类角色对订单的可见范围不同，联调过程中需要同步校验接口权限、前端路由守卫和页面按钮展示逻辑，避免出现“前端可见但后端无权”或“后端开放但前端入口缺失”的不一致情况。
\end{itemize}

\section{当前阶段状态}

\begin{itemize}[leftmargin=2em]
  \item 代码仓库已经形成较完整的前后端分离工程骨架，核心页面、接口与数据模型能够相互对应；
  \item 用户下单、代练申请、管理员审核、代练接单与订单状态推进等主流程已经在代码层面打通；
  \item 数据库结构已从基础用户/订单模型扩展到代练申请与审核机制，具备继续演进的基础；
  \item 当前版本已经具备课程中期演示条件，但支付结算、即时通讯、系统化审计与自动化测试仍属于后续补强内容。
\end{itemize}

\section{下一阶段建议}

\begin{itemize}[leftmargin=2em]
  \item 补充自动化测试，重点覆盖认证流程、订单状态流转、权限边界与并发接单场景；
  \item 强化上传安全能力，例如文件类型白名单、大小限制与内容校验，避免恶意文件进入系统；
  \item 继续完善密钥轮换与历史加密数据迁移策略，降低长期运维风险；
  \item 为管理员侧补充统计视图、操作审计与异常处理记录，提升后台治理能力；
  \item 在时间允许的前提下，推进即时通讯、支付结算或更稳健的 AI 兜底策略，使系统更接近真实业务场景。
\end{itemize}

\section*{附：演示材料说明}
本报告重点说明代码实现与阶段性工程进展。若配合演示视频或现场讲解，建议按照“用户注册登录 $\rightarrow$ AI 辅助下单 $\rightarrow$ 提交代练申请 $\rightarrow$ 管理员审核 $\rightarrow$ 代练接单与完单 $\rightarrow$ 异常订单处理”的顺序展开，以更清晰地呈现本阶段已经完成的业务闭环。

\end{document}

